%----------------------------------------------------------------------------------------
% Tailored CV — Computational Immunology PhD, UK Erlangen (Job 12900)
%----------------------------------------------------------------------------------------

\documentclass[10pt,a4paper,sans]{moderncv}

\usepackage[utf8]{inputenc}
\usepackage[T1]{fontenc}

\moderncvstyle{classic}
\moderncvcolor{blue}

\usepackage[scale=0.78]{geometry}
\usepackage{ragged2e}

%----------------------------------------------------------------------------------------
%	PERSONAL INFORMATION
%----------------------------------------------------------------------------------------

\firstname{Kundai Farai}
\familyname{Sachikonye}

\address{Auf der Lichtung 19}{82178 Puchheim, Germany}
\mobile{(+49) 1626386344}
\email{kundai.sachikonye@bitspark.com}
\homepage{github.com/fullscreen-triangle}{github.com/fullscreen-triangle}

%----------------------------------------------------------------------------------------

\begin{document}

\makecvtitle

%----------------------------------------------------------------------------------------
%	EDUCATION
%----------------------------------------------------------------------------------------

\section{Education}

\cventry{2019--2021}{PhD candidate, Computational Lipidomics}{Technische Universität München}{}{}{
  Doctoral research in mass spectrometry-based lipidomics, multi-omics statistical modelling,
  and scientific computing. Departed programme to pursue independent computational research.
}

\cventry{2018--2019}{MSc Computational Biology}{Universität Tübingen}{}{}{}

\cventry{2014--2017}{MSc Pharmaceutical Biotechnology}{HAW Hamburg}{}{}{}

\cventry{2011--2014}{BSc Biochemistry and Cell Biology}{Jacobs University Bremen}{}{}{}

%----------------------------------------------------------------------------------------
%	RESEARCH OUTPUT
%----------------------------------------------------------------------------------------

\section{Research Output}

\cvitem{2025}{
  \textbf{On the Thermodynamic Consequences of Categorical Completion on Biological
  Systems: Specificity through VDJ Recombination Categorical Filter Cascades.}
  \newline
  Proposes that adaptive immunity operates through categorical equivalence class filtering.
  MHC binding affinity correlates with parent protein categorical richness
  ($r=0.73$, $p<10^{-12}$, $n=2{,}847$ IEDB peptides). Adding the richness term to
  NetMHCpan improves AUC from 0.68 to 0.81. Framework predicts autoimmune targets
  (collagen, myelin, insulin) from protein conformational entropy and isoform diversity.
  \textit{Manuscript.}
}

\cvitem{2025}{
  \textbf{On Disease Epistemology in Blind Receiver: Cellular Ignorance, Group Blindness,
  and Self-Consistency as the Unique Basis of Disease Detection.}
  \newline
  Proves via two axioms that template-based diagnosis is epistemologically impossible at
  every biological scale. Disease is formalised as loop-holonomy defect in the biochemical
  circuit ($H_\ell \neq \mathrm{Id}$); holonomy test achieves AUC$=1.00$ vs 0.896 for
  template comparison. Therapy is sparse $\ell_1$ LP for loop-closure restoration.
  25/25 numerical validation experiments pass.
  \textit{Manuscript. AIMe Registry / TUM School of Life Sciences Weihenstephan.}
}

\cvitem{2025}{
  \textbf{Syndrome: Categorical Resolution of Disease Trajectories.}
  \newline
  Mathematical framework for disease trajectory computation through partition geometry.
  Eight cellular oscillator classes (protein, enzyme, channel, membrane, ATP, genetic
  expression, calcium, circadian) each yield a coherence index $\eta_i\in[0,1]$;
  disease vector $\mathbf{D}\in[0,1]^8$ encodes pathological signatures with
  $\mathcal{O}(N\log N)$ trajectory resolution.
  \textit{Manuscript.}
}

\cvitem{2025}{
  \textbf{Viral Infection as Host-State-Dependent Categorical Resonance.}
  \newline
  Derives infection probability as a categorical overlap integral from first principles
  (Maxwell electrostatics, quantum chemistry). Proves the Oscillatory Incompleteness
  Theorem: viral particles cannot orchestrate host machinery.
  \textit{Manuscript.}
}

%----------------------------------------------------------------------------------------
%	RESEARCH SOFTWARE
%----------------------------------------------------------------------------------------

\section{Research Software}

\cvitem{lavoisier}{
  Mass-spectrometry metabolomics analysis framework. Dual numerical and computer-vision
  pipeline; 98.9\% feature extraction accuracy against NIST glycan reference library;
  neural networks (PyTorch), distributed computing (Ray), supports $>$100\,GB mzML
  datasets via memory-mapped I/O.
  \href{https://github.com/fullscreen-triangle/lavoisier}{github.com/fullscreen-triangle/lavoisier}
}

\cvitem{gospel}{
  Whole-genome variant analysis and RNA-seq integration framework. Processes million-scale
  VCF files at $10^6$ variants/second; Bayesian network inference, Gaussian mixture models,
  variational Bayes. Validated against 1000 Genomes Phase 3, gnomAD v3.1.2, UK Biobank.
  \href{https://github.com/fullscreen-triangle/gospel}{github.com/fullscreen-triangle/gospel}
}

\cvitem{helicopter}{
  Multi-modal microscopy image analysis. Combines fluorescence, brightfield, phase-contrast,
  and spectral modalities through sequential exclusion to resolve structural ambiguity
  beyond single-modality optical limits.
  \href{https://github.com/fullscreen-triangle/helicopter}{github.com/fullscreen-triangle/helicopter}
}

\cvitem{syndrome}{
  Implementation of the categorical disease trajectory framework. Oscillator coherence
  computation, disease vector representation, trajectory resolution via constraint
  satisfaction. Python/Rust.
  \href{https://github.com/fullscreen-triangle/syndrome}{github.com/fullscreen-triangle/syndrome}
}

%----------------------------------------------------------------------------------------
%	RESEARCH EXPERIENCE
%----------------------------------------------------------------------------------------

\section{Research Experience}

\cventry{2019--2021}{Computational Lipidomics}{Bayerisches Forschungsinstitut}{Munich}{}{
  Mass spectrometry-based lipidomics pipeline: peak detection (Java), ion database
  annotation (PHP/Python), feature engineering, statistical modelling, and front-end
  visualisation of multi-omics datasets. Federated learning architecture for distributed
  analysis across clinical sites.
}

\cventry{2017--2018}{Glycomics and Proteomics}{Universitätsklinikum Hamburg-Eppendorf}{}{Python}{
  Automated container-based pipeline for standardised proteomics and glycomics
  characterisation of LC-MS/MS and MALDI-TOF data. High-throughput sample processing
  and statistical QC.
}

\cventry{2019}{Bioprocess Optimisation}{Fraunhofer Institut IGB}{Stuttgart}{Python, C}{
  Real-time process monitoring and data acquisition for industrial algae photobioreactors.
  Dynamic and linear programming methods for production optimisation of therapeutic lipids.
}

\cventry{2013}{Biomarker Discovery}{University Medical Center Groningen}{}{Python}{
  UPLC-MS/MS shotgun proteomics protocol optimisation for cervical cancer biomarker
  discovery in cell plasma and serum.
}

\cventry{2013}{Pharmacokinetics}{Leibniz Lungenzentrum}{}{Data Analysis}{
  Quantification of matrix suppression effects in anti-TB drug efficacy measurement
  in CHO lung cells using Selected Reaction Monitoring LC-MS.
}

\cventry{2012}{Metabolomics}{University of Zimbabwe}{}{Data Analysis}{
  Development and optimisation of a plant metabolite characterisation pipeline
  using LC and UV-spectrometry.
}

%----------------------------------------------------------------------------------------
%	TECHNICAL SKILLS
%----------------------------------------------------------------------------------------

\section{Technical Skills}

\cvitem{Languages}{Python, R, Rust, TypeScript/JavaScript, Java, C, SQL, Bash}

\cvitem{Bioinformatics}{
  Single-cell RNA-seq analysis, variant calling (VCF), mass spectrometry data processing
  (mzML, mzXML, FASTA, BAM), sequence alignment, pathway analysis (KEGG, Reactome),
  protein structure analysis (RDKit, GROMACS)
}

\cvitem{Statistics / ML}{
  Bayesian networks, Gaussian mixture models, variational Bayes, Gaussian processes,
  gradient boosting, graph neural networks, LSTMs, Transformers, PyTorch, scikit-learn,
  NumPy, SciPy, Pandas
}

\cvitem{Omics Platforms}{
  LC-MS/MS (QTOF, Quadrupole, Ion trap), MALDI-TOF, TIMS-PASEF, Native-MS,
  UPLC-MS/MS, qPCR, NGS, DLS, Fluorescence and Atomic Force Microscopy, SPR
}

\cvitem{Infrastructure}{
  Docker, distributed computing (Ray, Dask), HPC (Azure, AWS Lambda), Git,
  PostgreSQL, Neo4j, Jupyter
}

\cvitem{Data Formats}{
  mzML, mzXML, FASTA, BAM, VCF, HDF5, SBML, AniML, NeuroML, JSON
}

%----------------------------------------------------------------------------------------
%	LANGUAGES
%----------------------------------------------------------------------------------------

\section{Languages}

\cvitemwithcomment{English}{Native / Fluent}{}
\cvitemwithcomment{German}{Conversational}{Living and working in Germany since 2011}

%----------------------------------------------------------------------------------------

\renewcommand{\listitemsymbol}{-~}

\end{document}
